% !TEX root = AImath.v4.tex
\chapter*{ 前言}
%\markboth{前言}{前言}

\addcontentsline{toc}{chapter}{前言}

\begin{introduction}

\item 现在看起来的“显而易见”，在当年却是“困难重重”。这句话深刻地概括了人工智能发展的历史轨迹，也揭示了人类智慧演进的根本规律。今天我们习以为常的技术成就，在其诞生之初都曾面临着巨大的挑战和质疑。

\end{introduction}

\section{从“困难重重”到“显而易见”：智能发展的历史辩证法}

现在看起来的“显而易见”，在当年却是“困难重重”。这句话，也许是理解人工智能发展历程的最好钥匙。今天我们轻松使用语音助手、个性化推荐、智能搜索摘要、实时翻译，甚至让大模型生成文本、代码与图像，驱动办公 Copilot 与多模态对话——很难想象这些能力背后经历过怎样的艰难探索。每一个今天看来理所当然的功能，都是无数研究者在黑暗中摸索、在失败中坚持、在质疑中前行的结果。

人类历史上，每一次重大的技术突破都遵循着相似的轨迹：从被认为不可能，到被视为困难，再到变得显而易见，最终成为日常生活的一部分。这种认知转变的背后，是科学与工程共同推动的深层逻辑。

在古代，人们认为飞行是神的特权，莱特兄弟的尝试曾被视为“杂技”；几十年后，航空成为日常。人工智能亦如此：从图灵的“机器能思考吗？”，到深蓝与AlphaGo，再到通用大模型的涌现——不可能逐步变为显然。

在神话传统中，人造“智能”的想象从未缺席——塔罗斯与哥连都在提醒我们：人类对智能的创造有着恒久的渴望。要理解这种“从困难到显而易见”的逻辑，最清晰的路径，是回到数学与计算的演化主线。

\section{数学史：智能追求的发展主线}

本质上，人类发展史就是一部智能简史。从最早的数概念，到算术、代数、几何、抽象代数、拓扑、群论等日益复杂的分支，数学史的主线背后，是对更强智能的追求。

从最初的计数开始，人类就在不断追求计算更高效、更自动化、更智能。算盘让计算摆脱手指；对数表使复杂运算变得简单；机械计算器标志自动化的开端。每一步进展，都是在拓展智能的边界。从人脑和手算发展到一定程度，对更高效计算效率的追求必然导致计算机的发明，进而催生人工智能——这不是偶然，而是历史发展的必然逻辑。

“语言是思想的边界，技术是思想的实现。”新概念的诞生拓展思维边界，抽象理论化为算法与程序，思想由此获得改变世界的力量。深度学习的历程即是典型：从人工神经元到反向传播，再到卷积网络的突破，理论与实践相互促进，推动领域飞跃。

数学在其中扮演了底座角色：线性代数提供计算框架，概率论奠定学习理论，微积分使梯度优化成为可能，信息论指导数据压缩与特征提取。这些工具也经历了从“困难重重”到“显而易见”的转变——微积分自争议而至基础课程，概率论自“数学游戏”而成科学支柱。

欧几里得几何的公理化、牛顿微积分的发明、非欧几何的发现、集合论的建立、哥德尔不完备定理的证明——这些“宿命时刻”如星光点亮历史，照亮了智能进化的道路。这条主线直接推动了智能层次的跃迁——从模仿到创造。

这条主线也具体映射到人工智能能力的形成：优化与凸分析支撑支持向量机与深度网络的训练，图论与组合推动知识图谱与图神经网络，数值方法与稳定性理论保障大规模训练的可行性，几何与群论解释视觉中的不变性与等变性。由抽象到算法，由算法到系统，数学不断把“不可为”改写为“可为”。

\section{智能的层次：从模仿到创造}

智能不仅是计算能力，更是对复杂情境的拿捏与适度响应。“过犹不及”正是智能的外在表现之一。

最初的人工智能依赖规则与逻辑推理（符号主义），在复杂现实面前力有未逮；随后机器学习兴起，依靠数据统计与模式识别，能学能泛化却仍以模仿为主；深度学习通过多层非线性映射，展现出某种“理解”，能处理语言、图像与跨模态生成，但本质上仍以统计关联为核心。

从规则到学习、从统计到生成，智能的跃迁不是凭空发生；它背后站着一整套可操作的数学框架。要让这种跃迁站稳脚跟，还需要它的底座——数学。

真正的挑战在于让机器具备创造性思维与抽象推理能力。这一挑战在今日呈现为“工具使用与代理化”的能力：大模型结合检索增强、函数调用、程序生成与规划执行，开始从模仿走向协作与解决问题；但要让这种能力可控、可证、可复现，仍需符号推理与神经表征的更紧密融合。我们将在后文讨论“统计—符号”混合范式如何为此提供可能路径。

\section{本书的目标与挑战}

本书试图通过数学的视角，揭示人工智能发展的内在逻辑。我们将探讨那些看似“显而易见”的技术背后的数学原理，回顾那些曾经“困难重重”的问题如何被一步步解决。

人工智能高度跨学科，涉及数学、计算机科学、认知科学、神经科学、哲学等。要在有限篇幅内既保持严谨，又保证可读性，需要在深度与广度之间取得平衡。

写作方式上，我们采用“公式—直觉—应用—限制—展望”的五段式展开，并以“案例—数学—实现”三层结构贯穿，确保读者既能把握数学内核，又能落到工程实践。

我们的目标不仅是传授知识，更是启发思考：今天的“显而易见”，来源于昨日的“困难重重”；明天的突破，或许就蕴藏在今天看似不可能的想法之中。每一次低谷都被视为“AI 冬天”，每一次复兴都被视为“奇迹”。

在这个快速演进的时代，我们既要敬畏技术，也要清醒面对挑战。数学作为人工智能的基石，将继续指引我们探索智能的奥秘，推动文明迈向更高层次。

让我们带着这样的认识，开始这段探索人工智能数学基础的旅程。在这个旅程中，我们将见证“困难重重”如何转化为“显而易见”，体验数学之美如何照亮智能之路。

\section{致读者}

这本书不是传统教科书，也不是纯粹的学术专著。它更像是一场思想漫游——一次关于智能本质的哲学思辨，一段追溯数学与人工智能关系的历史探索。

我们将从基础数学概念出发，逐步深入到人工智能的核心算法。每一个公式背后都有思想；每一个算法背后都有故事。我们不仅要理解“是什么”“怎么做”，更要思考“为什么”“意味着什么”。

正如史蒂夫·乔布斯所言，连接点只能在回望时发生——当你一路走来，再回看这些“点”，你会发现它们自有逻辑地连成了线。愿这本书为你提供有意义的“点”，让你在未来的某刻完成自己的连接。

技术会过时，思想永恒；工具会更新，智慧长存。愿这本书成为你探索智能奥秘的指南针，在“困难重重”的学习路上点亮一盏明灯，最终让那些看似高深的理论变得“显而易见”。

\newpage